
\documentclass[12pt]{article}
\usepackage[utf8]{inputenc}
\usepackage[portuguese]{babel}
\usepackage{booktabs}
\usepackage{array}
\usepackage{geometry}
\usepackage{amsmath}
\usepackage{amsfonts}
\usepackage{longtable}

\geometry{a4paper, margin=1.5cm}

\title{Solução do Problema Sc50teste - Solver HiGHS}
\author{Análise Computacional}
\date{\today}

\begin{document}

\maketitle

\section{Informações do Problema}

\textbf{Informações do Problema:}
\begin{itemize}
\item Nome: Sc50teste
\item Número de Variáveis: 48
\item Número de Restrições: 50
\item Inviabilidade Primal: 0.000e+00
\item Inviabilidade Dual: 0.000e+00
\item Valor Primal: -6.458e+01
\item Valor Dual: -6.458e+01
\item Gap: 0.000e+00
\item Número de Iterações: 8
\end{itemize}


\section{Variáveis Primais e Custos Reduzidos (x > 0 e z > 0)}

\begin{longtable}{@{}cccc@{}}
\caption{Variáveis primais (x) e custos reduzidos (z) com valores > 0 do problema Sc50teste} \\
\toprule
\textbf{Coordenada x} & \textbf{Valor x (Primal)} & \textbf{Coordenada z} & \textbf{Valor z (Custo Reduzido)} \\
\midrule
\endfirsthead

\toprule
\textbf{Coordenada x} & \textbf{Valor x (Primal)} & \textbf{Coordenada z} & \textbf{Valor z (Custo Reduzido)} \\
\midrule
\endhead

\midrule \multicolumn{4}{r}{{Continua na próxima página}} \\ \midrule
\endfoot

\bottomrule
\endlastfoot
2 & 1.657e+01 & 9 & 7.706e-02 \\
3 & 6.458e+01 & 12 & 3.082e-02 \\
4 & 6.458e+01 &  &  \\
5 & 6.458e+01 &  &  \\
7 & 1.657e+01 &  &  \\
8 & 6.458e+01 &  &  \\
10 & 1.657e+01 &  &  \\
11 & 6.458e+01 &  &  \\
13 & 2.001e+01 &  &  \\
14 & 7.103e+01 &  &  \\
15 & 7.103e+01 &  &  \\
16 & 1.356e+02 &  &  \\
18 & 3.658e+01 &  &  \\
19 & 1.356e+02 &  &  \\
21 & 3.658e+01 &  &  \\
22 & 1.356e+02 &  &  \\
23 & 1.418e+01 &  &  \\
24 & 1.760e+01 &  &  \\
25 & 7.814e+01 &  &  \\
26 & 7.814e+01 &  &  \\
27 & 2.137e+02 &  &  \\
28 & 1.418e+01 &  &  \\
29 & 5.418e+01 &  &  \\
30 & 2.137e+02 &  &  \\
31 & 1.418e+01 &  &  \\
32 & 5.418e+01 &  &  \\
33 & 2.137e+02 &  &  \\
34 & 1.842e+01 &  &  \\
35 & 1.842e+01 &  &  \\
36 & 8.595e+01 &  &  \\
37 & 8.595e+01 &  &  \\
38 & 2.997e+02 &  &  \\
39 & 3.260e+01 &  &  \\
40 & 7.260e+01 &  &  \\
41 & 2.997e+02 &  &  \\
42 & 3.260e+01 &  &  \\
43 & 7.260e+01 &  &  \\
44 & 2.997e+02 &  &  \\
45 & 2.026e+01 &  &  \\
46 & 2.026e+01 &  &  \\
47 & 9.454e+01 &  &  \\
48 & 9.454e+01 &  &  \\

\bottomrule
\end{longtable}


\section{Variáveis Duais (Multiplicadores de Lagrange - Valores > 0)}

\begin{longtable}{@{}cc@{}}
\caption{Variáveis duais (y) com valores > 0 do problema Sc50teste} \\
\toprule
\textbf{Coordenada y} & \textbf{Valor y (Dual)} \\
\midrule
\endfirsthead

\toprule
\textbf{Coordenada y} & \textbf{Valor y (Dual)} \\
\midrule
\endhead

\midrule \multicolumn{2}{r}{{Continua na próxima página}} \\ \midrule
\endfoot

\bottomrule
\endlastfoot

\bottomrule
\end{longtable}


\end{document}
